\documentclass[10pt,a4paper]{amsart}
\usepackage[margin=19mm]{geometry}
\usepackage{amsmath,amssymb,mathtools}
\usepackage[scr=rsfs,cal=euler]{mathalfa}
\usepackage{xfrac}
\usepackage[unicode,hidelinks,bookmarks=false]{hyperref}
\newcommand{\ie}{i.e.\ }
\newcommand{\cf}{cf.\ }
\newcommand{\resp}{respectively\ }
\newcommand{\Subsection}{\subsection}
\providecommand{\nobreakdash}{\nobreak\hbox{-}}
\setlength{\emergencystretch}{2em}
\multlinegap0pt
\usepackage{bm}
\usepackage{bbm}
\usepackage{dsfont}
\usepackage{xspace}
\usepackage{cancel}
\usepackage{mathtools}
\usepackage{amsthm}
\makeatletter
\@ifundefined{theorem}{\newtheorem{theorem}{Theorem}}{}
\@ifundefined{lemma}{\newtheorem{lemma}[theorem]{Lemma}}{}
\makeatother
\def \QR{\mathbb{Q}_\mathbb{R}}
\def \R{{\mathbb R}}
\def \i{\textit{\textbf{i}}}
\def \j{\textit{\textbf{j}}}
\def \k{\textit{\textbf{k}}}
\newcommand\norm[1]{\left\lVert#1\right\rVert}
\title[Solution Methods and Perturbation Analysis for the Equality-]{Solution Methods and Perturbation Analysis for the Equality-Constrained Total Least Squares Problem over Reduced Biquaternions}
\author{Neha Bhadala}
\date{Source: arXiv:2506.18319v1 (CC BY 4.0)}
\begin{document}
\maketitle

\noindent\emph{Source and licence.} Solution Methods and Perturbation Analysis for the Equality-Constrained Total Least Squares Problem over Reduced Biquaternions. Version arXiv:2506.18319v1 (CC BY 4.0), distributed under the Creative Commons Attribution 4.0 International licence, as recorded in the arXiv licence metadata for this exact version and shown on the abstract page https://arxiv.org/abs/2506.18319v1. Full rights notice: licenses/arxiv-2506.18319-CC-BY-4.0.txt.\par\smallskip

\noindent\textit{Editorial scope.} Counted unit: the Theorem whose source label is \texttt{thm3.1} in the article (source0.tex lines 231--244, source label \texttt{thm3.1}), delivered together with the article's own complete original proof of it (source0.tex lines 245--343, one \texttt{proof} environment of 99 lines). No mathematics is written, repaired or re-derived: statement and proof are the article's own text line for line. Required context retained uncounted: eq2.2 (lines 138--148); eq2.3 (lines 171--173); lem2 (lines 191--198); eq3.3 (lines 217--219); eq3.4 (lines 224--226). Every definition, display and label of those lines is retained unchanged. Omitted from this package and not redistributed: 1--137, 149--170, 174--190, 199--216, 220--223, 227--230, 344--1286. Nothing inside a retained excerpt is omitted or altered: every retained body is delivered exactly as the article's lines are written, so no omission receipt of any kind accompanies this package. Citations: the retained excerpts carry no citation command at all, so no bibliography entry of the article is transcribed and no reference is redistributed. Reference adaptation: none was needed and none was made; every reference inside the delivered unit resolves inside it, so no unresolved reference or citation is produced. Printed slips kept literal: no printed word, symbol, hypothesis or number of the counted statement or of its proof was altered. Layout and class: the article's own class is \texttt{elsarticle} with packages including amsmath, amsthm, amsfonts, amssymb, xcolor, algorithm, multirow, algpseudocode, blkarray, booktabs, mathrsfs, sectsty, hyperref, subfigure; the wrapper is the lane's pinned \texttt{amsart} editorial wrapper at 10pt on A4, which restates the theorem-like environments the retained text needs and adds the article's own macro definitions that the retained text needs (\texttt{\textbackslash QR}, \texttt{\textbackslash R}, \texttt{\textbackslash i}, \texttt{\textbackslash j}, \texttt{\textbackslash k}, \texttt{\textbackslash norm}), with extra wrapper packages (bm, bbm, dsfont, xspace, cancel, mathtools). The delivered numbering is the unit's own. Assets: no retained span contains a figure environment, a graphics-inclusion command, a PDF-page inclusion, a table or a TikZ picture; no image file is copied into this package, the package declares no asset dependency and no image file is redistributed with it. Licence: version arXiv:2506.18319v1 of this article is distributed under the Creative Commons Attribution 4.0 International licence, as recorded in the arXiv licence metadata for this exact version and shown on the abstract page https://arxiv.org/abs/2506.18319v1. A full rights notice is delivered at licenses/arxiv-2506.18319-CC-BY-4.0.txt. Characters: every retained excerpt is pure ASCII, so the delivered engine, XeLaTeX with its default Latin Modern text font, reports no missing character.
\begin{equation}\label{eq2.2}
	P^R= \begin{bmatrix}
		P_0  &  -P_1  &  P_2  &  -P_3 \\
		P_1   &   P_0  &  P_3  &   P_2  \\
		P_2   &  -P_3 & P_0  &  -P_1  \\
		P_3  &  P_2  &  P_1  &  P_0 
	\end{bmatrix}, ~ P^C= \begin{bmatrix}
		R_1 & R_2 \\
		R_2 & R_1 
	\end{bmatrix}.
\end{equation}

\begin{equation}\label{eq2.3}
	P^R= [P_c^R, \mathcal{K}_m P_c^R, \mathcal{L}_m P_c^R, \mathcal{M}_m P_c^R].  
\end{equation}

\begin{lemma}\label{lem2}
	Let $P= P_0+ P_1\i+ P_2\j+ P_3\k \in \QR^{m \times n}$ and $Q= Q_0+ Q_1\i+ Q_2\j+ Q_3\k \in \QR^{m \times d}$. Then, the following results hold.
	\begin{enumerate}
		 \item[\rm (i)] $\|[P, Q]\|_F= \frac{1}{2}\|[P, Q]^R\|_F$.
		 \item[\rm (ii)] $\|[P, Q]^R\|_F= \|[P^R, Q^R]\|_F$.
		\item[\rm (iii)] $\|[P^R, Q^R]\|_F= 2\|[P_c^R, Q_c^R]\|_F$.
	\end{enumerate}
\end{lemma}

\begin{equation}\label{eq3.3}
	\min_{X,\bar{E},\bar{F}}\norm{[\bar{E}, \bar{F}]}_F ~ \textrm{subject to} ~  (A+\bar{E})X = B+\bar{F}, ~ CX=D. 
\end{equation} 

\begin{equation}\label{eq3.4}
	\min_{X,\widetilde{E}, \widetilde{F}} \norm{[\widetilde{E}, \widetilde{F}]}_F ~ \textrm{subject to} ~ (A^R_c+\widetilde{E})X=B^R_c+ \widetilde{F}, ~ C^R_cX=D^R_c.
\end{equation}

\begin{theorem}\label{thm3.1}
	A matrix $X \in \R^{n \times d}$ is a real RBTLSE solution of \eqref{eq3.3} if and only if it is a real TLSE solution of \eqref{eq3.4}. Furthermore, let $X$ be a real TLSE solution with corresponding minimizing perturbations $\widetilde{E}$ and $\widetilde{F}$ in equation \eqref{eq3.4} given by
	\begin{equation*}
		\widetilde{E} = \begin{bmatrix} E_0^T & E_1^T & E_2^T & E_3^T \end{bmatrix}^T \in \R^{4m \times n}, ~
		\widetilde{F} = \begin{bmatrix} F_0^T & F_1^T & F_2^T & F_3^T \end{bmatrix}^T \in \R^{4m \times d},
	\end{equation*}
where $E_t \in \R^{m \times n}$ and $F_t \in \R^{m \times d}$ for $t=0,1,2,3$.

Then the corresponding minimizing perturbations $\bar{E}$ and $\bar{F}$ in equation \eqref{eq3.3} are given by
\begin{equation*}
	\bar{E} = E_0 + E_1\i + E_2\j + E_3\k \in \QR^{m \times n}, ~
	\bar{F} = F_0 + F_1\i + F_2\j + F_3\k \in \QR^{m \times d}.
\end{equation*}
\end{theorem}

\begin{proof}
	Let $X \in \R^{n \times d}$ be a real TLSE solution. Then, there exist perturbation matrices $\widetilde{E}$ and $\widetilde{F}$ satisfying
	\begin{equation*}
		\norm{[\widetilde{E}, \widetilde{F}]}_F= \textrm{min}, ~	(A^R_c+ \widetilde{E})X= B^R_c+ \widetilde{F}, ~ C^R_cX=D^R_c.
	\end{equation*}
	Using these equations, we get
	\begin{align}\label{eq3.5}
		\left[(A^R_c+ \widetilde{E}), \mathcal{K}_m (A^R_c+ \widetilde{E}), \mathcal{L}_m (A^R_c+ \widetilde{E}), \mathcal{M}_m (A^R_c+ \widetilde{E})\right]
		\begin{bmatrix}
			X & 0 & 0 & 0 \\
			0 & X & 0 & 0 \\
			0 & 0 & X & 0 \\
			0 & 0 & 0 & X
		\end{bmatrix} \nonumber &\\= \left[(B^R_c+ \widetilde{F}), \mathcal{K}_m (B^R_c+ \widetilde{F}), \mathcal{L}_m (B^R_c+ \widetilde{F}), \mathcal{M}_m (B^R_c+ \widetilde{F}) \right],
	\end{align}
and
	\begin{equation}\label{eq3.51}
		\left[C^R_c, \mathcal{K}_p C^R_c, \mathcal{L}_p C^R_c, \mathcal{M}_p C^R_c\right]	\begin{bmatrix}
			X & 0 & 0 & 0 \\
			0 & X & 0 & 0 \\
			0 & 0 & X & 0 \\
			0 & 0 & 0 & X
		\end{bmatrix} = \left[D^R_c, \mathcal{K}_p D^R_c, \mathcal{L}_p D^R_c, \mathcal{M}_p D^R_c\right].
	\end{equation}
Next, express the perturbed matrices as
	\begin{equation}\label{eq3.6}
		A^R_c + \widetilde{E}= 
		\begin{bmatrix}
			A_{0}+ E_{0} \\
			A_{1}+ E_{1} \\
			A_{2}+ E_{2} \\
			A_{3}+ E_{3} 
		\end{bmatrix}, ~ B^R_c+\widetilde{F}= 
		\begin{bmatrix}
			B_0+ F_0 \\
			B_1+ F_1  \\
			B_2+ F_2 \\
			B_3+ F_3
		\end{bmatrix}.
	\end{equation}
	Now define the RB matrices:
	\begin{eqnarray*}
		\bar{A}&:=& (A_{0}+E_{0})+ (A_{1}+E_{1})\i+ (A_{2}+E_{2})\j+ (A_{3}+E_{3})\k, \\
		\bar{B}&:=& (B_0+ F_0)+ (B_1+ F_1)\i+ (B_2+ F_2)\j+ (B_3+ F_3)\k, \\
		\bar{E}&:=& E_{0}+ E_{1}\i+ E_{2}\j+ E_{3}\k, ~ \bar{F}:= F_0+ F_1\i+ F_2\j+ F_3\k.
	\end{eqnarray*}
	Using equations \eqref{eq2.2} and \eqref{eq3.6}, we obtain
	\begin{equation}\label{eq3.61}
		\bar{A}_c^R= A^R_c+ \widetilde{E}, ~	\bar{B}_c^R= B^R_c+\widetilde{F}, ~ \bar{E}_c^R=  \widetilde{E}, ~ \textrm{and} ~ \bar{F}_c^R= \widetilde{F}.
	\end{equation}
Furthermore, using equations \eqref{eq2.3} and \eqref{eq3.61}, we can express $\bar{A}^R$, $\bar{B}^R$, and $X^R$ as
	\begin{eqnarray*}
		\bar{A}^R&=& \left[(A^R_c+ \widetilde{E}), \mathcal{K}_m(A^R_c+ \widetilde{E}), \mathcal{L}_m(A^R_c+ \widetilde{E}), \mathcal{M}_m(A^R_c+ \widetilde{E})\right], \\
		\bar{B}^R&=& \left[(B^R_c+\widetilde{F}), \mathcal{K}_m(B^R_c+\widetilde{F}), \mathcal{L}_m(B^R_c+\widetilde{F}), \mathcal{M}_m(B^R_c+\widetilde{F})\right],\\
		X^R&=& 
		\begin{bmatrix}
			X & 0 & 0 & 0 \\
			0 & X & 0 & 0 \\
			0 & 0 & X & 0 \\
			0 & 0 & 0 & X
		\end{bmatrix}.
	\end{eqnarray*}
	Thus, equation \eqref{eq3.5} can be rewritten as
	\begin{eqnarray}
		\bar{A}^R X^R&=& \bar{B}^R, \label{eq3.7}\\
		(\bar{A}X)^R&=& \bar{B}^R, \nonumber\\
		\bar{A}X&=& \bar{B}. \label{eq3.8}
	\end{eqnarray}
	The matrices $\bar{A}$ and $\bar{B}$ can be represented using $A$, $B$, $\bar{E}$, and $\bar{F}$ as follows
	\begin{eqnarray}
		\bar{A}&=& (A_{0}+ E_{0})+ (A_{1}+ E_{1})\i+ (A_{2}+ E_{2})\j+ (A_{3}+ E_{3})\k \nonumber \\
		&=& (A_0+ A_1\i+ A_2\j+ A_3\k)+ (E_0 +E_1\i+ E_2\j+ E_3\k) = A+ \bar{E}, \label{eq3.9}
	\end{eqnarray}
	and
	\begin{eqnarray}
		\bar{B}&=& (B_0+ F_0)+ (B_1+ F_1)\i+ (B_2+ F_2)\j+ (B_3+ F_3)\k \nonumber\\
		&=& (B_0+ B_1\i+ B_2\j+ B_3\k)+ (F_0 + F_1\i+ F_2\j+ F_3\k) = B+ \bar{F}. \label{eq3.10}
	\end{eqnarray}
	Using \eqref{eq3.9} and \eqref{eq3.10}, equation \eqref{eq3.8} is equivalent to
	\begin{equation} \label{eq3.11}
		(A+\bar{E})X= B+ \bar{F}.
	\end{equation}
	Also, using equation \eqref{eq2.3}, we have
	\begin{equation*}
		C^R = \left[C^R_c, \mathcal{K}_p C^R_c, \mathcal{L}_p C^R_c, \mathcal{M}_p C^R_c\right], ~ D^R=\left[D^R_c, \mathcal{K}_p D^R_c, \mathcal{L}_p D^R_c, \mathcal{M}_p D^R_c\right].
	\end{equation*}
Thus, equation \eqref{eq3.51} is equivalent to
	\begin{eqnarray}
		C^R X^R&=& D^R, \nonumber\\
		(CX)^R&=& D^R, \nonumber\\
		CX&=& D. \label{eq33.13}
	\end{eqnarray}
	Finally, using Lemma~\ref{lem2}, we relate the norms 
	\begin{equation}\label{eq3.12}
		\norm{[\bar{E}, \bar{F}]}_F= \frac{1}{2} \norm{[\bar{E}, \bar{F}]^R}_F= \frac{1}{2} \norm{[\bar{E}^R, \bar{F}^R]}_F =\norm{[\bar{E}^R_c, \bar{F}^R_c]}_F=  \norm{[\widetilde{E}, \widetilde{F}]}_F= \min.  
	\end{equation}
	Combining \eqref{eq3.11}, \eqref{eq33.13}, and \eqref{eq3.12}, we conclude that there exist matrices $\bar{E} \in \QR^{m \times n}$ and $\bar{F} \in \QR^{m \times d}$ such that
	$X  \in  \R^{n \times d}$ is a real RBTLSE solution, and vice versa.
\end{proof}

\end{document}
